\documentclass[12pt]{article}
\usepackage{amsmath, amssymb, enumitem, fancyhdr, graphicx, libertine, nth,
  pgfplots, tikz, wasysym, xcolor}
\pgfplotsset{compat=1.11}
\usetikzlibrary{decorations.pathreplacing}
\usetikzlibrary{patterns}
\usepackage[margin=1in]{geometry}
\usepackage[libertine]{newtxmath}
\pagestyle{fancy}
\definecolor{solution-color}{HTML}{000080}
\chead{\emph{EC201 -- Practice Final}}

% Define new topic command
\newcommand{\topic}[1]{\medskip\begin{center}\large\textsc{#1}\normalsize
  \end{center}}

% Define question counter:
\newcounter{question}[section]
\newcommand{\question}[1]{\refstepcounter{question}
  \subsubsection*{Question~\thequestion~-- #1}}

% Subquestions are enumerate with letters
\newenvironment{subquestions}{
  \begin{enumerate}[label=(\roman*)]}{\end{enumerate}}

\begin{document}
\question{Preferences (10 Points)}
\begin{itemize}
  \item Draw a pair of indifference curves that would violate transitivity.
    Explain why they violate transitivity.
  \item Draw an indifference curve that violates convexity. Explain why it
    violates convexity.
  \item Draw a pair of indifference curves that violate monotonicity. Explain
    why they violate monotonicity.
\end{itemize}
\question{Demand (10 Points)}
  Your utility function for goods 1 and 2 is given by:
  \[
    u\left( x_1,x_2 \right)= \log\left({x_1}\right) + x_2
  \]
\textbf{Note}: If $f\left( x \right)=\log\left( x \right)$, then
$\frac{df\left( x \right)}{dx}=\frac{1}{x}$.
\begin{subquestions}
\item Good 1 costs \$1 and good 2 costs \$2. You have \$10 to
  spend.  How much of good 1 will you buy and how much of good 2 will you buy?
\item For each good, are they normal or inferior?
\end{subquestions}
\question{Intertemporal Choice (10 Points)}
Your lifetime utility function is:
\[
u\left( c_1,c_2 \right)=c_1^{\frac{1}{2}}c_2^{\frac{1}{2}}
\]
The interest rate is 10\%. You earn income of \$100 in
period 1 and \$110 in period 2. How much will you save/borrow in period 1?

\question{Uncertainty (10 Points)}
Your total assets are worth \$30,000. Included in that \$30,000 are \$20,000
worth of valuable items in your apartment. There is a 10\% chance that you will be
burgled. The local insurance company offers to insure each dollar of stolen
items at 15\cent. If you fully insure, it will cost $\$20000\times 15\cent = \$3,000$. The remaining assets (the \$10,000) are safely locked in a
secure bank vault.

Let state 1 be the case where you are burgled and state 2 be the state where
you are not burgled. Call $c_1$ and $c_2$ your consumption in both states.
\begin{subquestions}
  \item Draw your budget contstraint for consumption across states where $c_1$
    is on the horizontal axis and $c_2$ is on the vertical axis. Assume you
    can't ``over insure'' (you can only buy a maximum of \$20,000-worth of
    insurance) and you can't negatively insure. You are, however, able to
    partially insure. Label the consumption path where you don't buy insurance
    and where you fully insure.  For each of these points, label how much you
    have to consume in each state of the world.
  \item What is the slope of the budget line?
  \item How much would an actuarially fair insurance policy cost?
  \item What would the slope of the budget line be if the insurance was
    actuarially fair? Would it be steeper or flatter?
\end{subquestions}

\question{Equilibrium, Firm Supply and Industry Supply (12 Points)}
The demand curve for a particular market is given by:
\[
  D\left( p \right)=880-20p
\]
There are 100 firms operating in this market each with a cost function $c\left(
q \right)=\frac{q^2}{4}$.
\begin{subquestions}
  \item What is the equilibrium price and quantity?
\end{subquestions}
Suddenly this good increases in popularity everyone is now willing to pay \$1
more for the good. The government decides now is a good time to introduce a
value tax of 25\% on the good.
\begin{subquestions}
  \setcounter{enumi}{1}
  \item What is the new equilibrium price and quantity as a result of the
    popularity increase and the tax? (Assume we're in the short run so that the
    number of firms does not change).
\end{subquestions}


\question{Monopoly (13 Points)}
Suppose market demand curve for a particular good is given by $D\left( p
\right)=70-\frac{1}{3}p$. There is only one firm selling this particular good.
The cost function for this firm is $c\left( q \right)=30q+3q^2$.
\begin{subquestions}
  \item What price and quantity does the monopolist choose? What is its profits?
  \item What is the elasticity of demand at the price and quantity the
    monopolist is operating at?
  \item If the monopolist could perfectly price discriminate, what profits
    could it achieve?
\end{subquestions}
\question{Game Theory (15 Points)}
\begin{subquestions}
  \item Consider the following game:
    \setlength{\arraycolsep}{0pt}
    \[
      \begin{array}{cccccc}
        &  & \quad L \quad & & \quad R \quad  &  \\ \cline{3-6} U \quad & \vline{}
        & \quad 2,4  \quad & \vline{} & \quad 1,3 \quad & \vline{}  \\ \cline{3-6}
        M \quad & \vline{} & \quad 4,2 \quad & \vline{} & \quad 0,-1 \quad &
        \vline{} \\ \cline{3-6} D \quad & \vline{} & \quad 3,1 \quad & \vline{} &
        \quad 5,3 \quad & \vline{} \\ \cline{3-6}
      \end{array}
    \]
    \begin{itemize}
      \item Does any player have any dominated strategies?
      \item Find all pure-strategy Nash equilibria of the game.
    \end{itemize}
    \item Consider the following extensive form game:
          \ \\
      \begin{center}
        \begin{tikzpicture}[scale=1.5,font=\footnotesize]
          \tikzstyle{solid node}=[circle,draw,inner sep=1.5,fill=black]
          \tikzstyle{hollow node}=[circle,draw,inner sep=1.5]
          \tikzstyle{level 1}=[level distance=15mm,sibling distance=3.5cm]
          \tikzstyle{level 2}=[level distance=15mm,sibling distance=1.5cm]
          \tikzstyle{level 3}=[level distance=15mm,sibling distance=1cm]
          \node(0)[hollow node,label=above:{$P1$}]{}
          child{node[solid node,label=above left:{$P2$}]{}
          child{node[label=below:{$(5,2)$}]{} edge from parent node[left]{$A$}}
          child{node[label=below:{$(3,-1)$}]{} edge from parent node[left]{$B$}}
          child{node[label=below:{$(0,3)$}]{} edge from parent node[right]{$C$}}
          edge from parent node[left,xshift=-5]{$L$}
          }
          child{node[solid node,label=above right:{$P2$}]{}
          child{node[label=below:{$(2,2)$}]{} edge from parent node[left]{$D$}}
          child{node[label=below:{$(4,1)$}]{} edge from parent node[right]{$E$}}
          edge from parent node[right,xshift=5]{$M$}
          }
          child[grow=right]{node{$(4,0)$} edge from parent node[above]{$R$}};
        \end{tikzpicture}
      \end{center}
      Find the subgame-perfect Nash equilibirium of the game.
    \end{subquestions}
\question{Oligopoly (20 Points)}
The inverse demand curve for a product is $p\left( q \right)=400-q$. Two firms
operate in this industry. Their marginal cost of production is constant at 100
and they have no fixed costs.

In this industry, the firms produce their output and then hand over all of
their output to a government distributor. The government distributor then
decides what price to charge such that all of the output will be sold. The
government distributor then gives the respective revenues back to the firms.
The government distributor does not charge a fee and has no costs.
\begin{subquestions}
  \item If the two firms compete by choosing quantities simultaneously, what
    output will each produce in equilibrium?  What will the price be? What will
    each firm's profits be?
  \item If one firm produces first and then publicly announces how much it
    produced before the other firm decides how much to produce, how much will
    each produce? What will the price be? What will their profits be?
  \item If both firms secretly decide together how much to produce in order to
    jointly maximize profits, how much will they produce? Assume here that if
    one firm cheats on the collusion quantity that the other firm will respond
    very aggresively, so neither firm will want to cheat on the collusion
    quantity. What will the price be? What will their profits be?
\end{subquestions}

\end{document}
